% Auto-generated by extract-marking.py -- do not edit by hand unless you know what to do.
% Works for BOTH \documentclass[...,type=solutions,...] and
% [...,type=marking,...] compiles, chosen at compile time via
% \OlympMarkPick / \OlympMarkBeginTable (see olympiad-marking.sty).
% \eqref{n} falls back to a plain '(n)' whenever the referenced
% equation isn't actually typeset in the current document.

\OlympMarkBeginTable{|p{0.862\linewidth}|>{\centering\arraybackslash}p{0.098\linewidth}|}{|p{0.655\linewidth}|>{\centering\arraybackslash}p{0.085\linewidth}|>{\centering\arraybackslash}p{0.06\linewidth}|>{\centering\arraybackslash}p{0.06\linewidth}|}
\hline
\OlympMarkPick{\multicolumn{1}{|c|}{\textbf{Мазмұны}} & \textbf{Ұпайлар} \\ \hline}{\multicolumn{1}{|c|}{\textbf{Мазмұны}} & \textbf{Ұпайлар} & \textbf{MP} & \textbf{A} \\ \hline}
\OlympMarkPick{Формула \eqref{1}: $\vec{v}_1(x,y,t) = \left(-\frac{v}{\sqrt{2}}, -\frac{v}{\sqrt{2}}\right) \cos\left(\omega t + \frac{\sqrt{2}\pi}{\lambda} x + \frac{\sqrt{2}\pi}{\lambda} y\right),$ & 0.4 \\ \hline}{Формула \eqref{1}: $\vec{v}_1(x,y,t) = \left(-\frac{v}{\sqrt{2}}, -\frac{v}{\sqrt{2}}\right) \cos\left(\omega t + \frac{\sqrt{2}\pi}{\lambda} x + \frac{\sqrt{2}\pi}{\lambda} y\right),$ & 0.4 & &  \\ \hline}
\OlympMarkPick{\quad\textit{\textbf{0.2 + 0.2 } вектордың дұрыс амплитудасы және фазаның координаттарға тәуелділігі үшін} &  \\ \hline}{\quad\textit{\textbf{0.2 + 0.2 } вектордың дұрыс амплитудасы және фазаның координаттарға тәуелділігі үшін} &  & &  \\ \hline}
\OlympMarkPick{Формула \eqref{2}: $\vec{v}_2(x,y,t) = \left(\frac{v}{\sqrt{2}}, -\frac{v}{\sqrt{2}}\right) \cos(\omega t - k_0 x + k_0 y),$ & 0.1 \\ \hline}{Формула \eqref{2}: $\vec{v}_2(x,y,t) = \left(\frac{v}{\sqrt{2}}, -\frac{v}{\sqrt{2}}\right) \cos(\omega t - k_0 x + k_0 y),$ & 0.1 & &  \\ \hline}
\OlympMarkPick{Формула \eqref{3}: $\vec{v}_3(x,y,t) = \left(-\frac{v}{\sqrt{2}}, \frac{v}{\sqrt{2}}\right) \cos(\omega t + k_0 x - k_0 y),$ & 0.1 \\ \hline}{Формула \eqref{3}: $\vec{v}_3(x,y,t) = \left(-\frac{v}{\sqrt{2}}, \frac{v}{\sqrt{2}}\right) \cos(\omega t + k_0 x - k_0 y),$ & 0.1 & &  \\ \hline}
\OlympMarkPick{Формула \eqref{4}: $\vec{v}_4(x,y,t) = \left(\frac{v}{\sqrt{2}}, \frac{v}{\sqrt{2}}\right) \cos(\omega t - k_0 x - k_0 y),$ & 0.1 \\ \hline}{Формула \eqref{4}: $\vec{v}_4(x,y,t) = \left(\frac{v}{\sqrt{2}}, \frac{v}{\sqrt{2}}\right) \cos(\omega t - k_0 x - k_0 y),$ & 0.1 & &  \\ \hline}
\OlympMarkPick{\quad\textit{Жоғарыдағы балдар \textbf{тек} вектордың дұрыс амплитудасы \textbf{және} фазасы үшін беріледі} &  \\ \hline}{\quad\textit{Жоғарыдағы балдар \textbf{тек} вектордың дұрыс амплитудасы \textbf{және} фазасы үшін беріледі} &  & &  \\ \hline}
\OlympMarkPick{Формула \eqref{5}: $v_y = 2\sqrt{2}v \cos(k_0 x) \sin(k_0 y) \sin(\omega t),$ , немесе жылдамдық компоненттерінің кем дегенде біреуі үшін баламалы қос тұрғын толқын теңдеуі үшін & 0.8 \\ \hline}{Формула \eqref{5}: $v_y = 2\sqrt{2}v \cos(k_0 x) \sin(k_0 y) \sin(\omega t),$ , немесе жылдамдық компоненттерінің кем дегенде біреуі үшін баламалы қос тұрғын толқын теңдеуі үшін & 0.8 & &  \\ \hline}
\OlympMarkPick{Формула \eqref{6}: $v_{\max} = \sqrt{8v^2} = 2\sqrt{2}v,$ & 0.5 \\ \hline}{Формула \eqref{6}: $v_{\max} = \sqrt{8v^2} = 2\sqrt{2}v,$ & 0.5 & &  \\ \hline}
\OlympMarkPick{Формула \eqref{7}: $\left( \frac{\lambda}{2 \sqrt{2}} \cdot n, \frac{\lambda}{2 \sqrt{2}} \cdot m \right) \ \ \text{мұнда \(n,m\) сандарының тек біреуі тақ}.$ & 1.0 \\ \hline}{Формула \eqref{7}: $\left( \frac{\lambda}{2 \sqrt{2}} \cdot n, \frac{\lambda}{2 \sqrt{2}} \cdot m \right) \ \ \text{мұнда \(n,m\) сандарының тек біреуі тақ}.$ & 1.0 & &  \\ \hline}
\OlympMarkPick{\quad\textit{\textbf{1.0} максимум координаттары үшін толық дұрыс және негізделген жауап үшін} &  \\ \hline}{\quad\textit{\textbf{1.0} максимум координаттары үшін толық дұрыс және негізделген жауап үшін} &  & &  \\ \hline}
\OlympMarkPick{\quad\textit{\textbf{0.4} жартылай дұрыс жауап үшін, егер максимум нүктелері тікбұрышты \textbf{шексіз торда} болатыны көрініп тұрса} &  \\ \hline}{\quad\textit{\textbf{0.4} жартылай дұрыс жауап үшін, егер максимум нүктелері тікбұрышты \textbf{шексіз торда} болатыны көрініп тұрса} &  & &  \\ \hline}
\OlympMarkPick{\textbf{Барлығы} & \textbf{3.0} \\ \hline}{\textbf{Барлығы} & \textbf{3.0} &  &  \\ \hline}
\end{tabularx}
